% !TeX program = pdflatex
% ==========================================================================
%  HPDC-MaOEA -- two-column manuscript draft (Elsevier elsarticle, 5p layout)
%
%  The manuscript is organised around two contributions: hybrid PBI quality
%  grouping and dual-mode candidate selection. Established components are
%  cited where they enter the framework.
%
%  Layout notes:
%    * elsarticle option "5p" is Elsevier's final two-column journal layout
%      (textwidth 522pt, columnsep 18pt => 252pt columns) and it loads
%      fleqn.clo, so displayed equations are flush left.
%    * Wide displays are broken with split/multline and wide floats use the
%      starred environments (table*, algorithm*) to span both columns.
%      Line breaking and float placement are adapted to the column width; the
%      wording and the mathematics are not affected by it.
%    * For the double-spaced single-column referee version, swap the class
%      options for [preprint,review,12pt].
% ==========================================================================
\documentclass[final,5p,times,twocolumn]{elsarticle}

\usepackage{amsmath,amssymb,amsthm}
\usepackage{bm}
\usepackage{booktabs}
\usepackage{algorithm}
\usepackage{algpseudocode}
\usepackage{enumitem}
\usepackage{xcolor}
\usepackage[hidelinks]{hyperref}

% graphicx is already loaded by elsarticle.cls

\newtheorem{proposition}{Proposition}
\newtheorem{remark}{Remark}

% Narrow columns: allow slightly looser lines before an overfull box.
\tolerance=1200
\emergencystretch=1em

% Pseudo-code is set one size down so that the wider \State lines fit a column.
\usepackage{etoolbox}
\AtBeginEnvironment{algorithmic}{\footnotesize}
\algrenewcommand{\algorithmicrequire}{\textbf{Input:}}
\algrenewcommand{\algorithmicensure}{\textbf{Output:}}

\newcommand{\vf}{\mathbf{f}}
\newcommand{\vx}{\mathbf{x}}
\newcommand{\vv}{\mathbf{v}}
\newcommand{\vw}{\mathbf{w}}
\newcommand{\vr}{\mathbf{r}}
\newcommand{\vz}{\mathbf{z}}
\newcommand{\TODO}[1]{\textcolor{red}{[TODO: #1]}}

%% Uncomment and fill in for a preprint that names its target journal.
%% \journal{Information Sciences}

\begin{document}

\begin{frontmatter}

\title{Hybrid PBI-Based Quality Grouping and Dual-Mode Candidate Selection
for Expensive Many-Objective Optimization}

%% Author block. Replace the placeholders and add \corref/\ead as needed:
%% \author[inst1]{Given Name Surname\corref{cor1}}
%% \ead{name@example.edu}
%% \cortext[cor1]{Corresponding author.}
\author{\TODO{author list}}

\affiliation{organization={Department, University},
            addressline={Street},
            city={City},
            postcode={000000},
            state={Province},
            country={Country}}

\begin{abstract}
\TODO{Use the challenge--contribution structure. State the two concrete gaps:
binary grouping alone provides no within-label ordering for positive-group selection,
and a fixed relation-score rule lets the score distribution determine how many
candidates are evaluated. Introduce hybrid PBI-based quality grouping and
dual-mode candidate selection as two coordinated design components, addressing
training-group construction and expensive-evaluation selection. Lead the results
with the strongest verified algorithm-level IGD comparison in the tested
many-objective settings. Then give one concise mechanism finding about group
composition or candidate-batch utility. Keep detailed overlap statistics and
componentwise tests in the experimental section.}
\end{abstract}

\begin{keyword}
expensive many-objective optimization \sep surrogate-assisted evolutionary
algorithm \sep relation learning \sep penalty-based boundary intersection \sep
candidate selection
\end{keyword}

\end{frontmatter}

\section{Introduction}
\label{sec:intro}

\TODO{Use a four-move opening: expensive many-objective relation learning; the
loss of within-label ordering when forming the positive group and the dependence of
fixed-cutoff batch size on the relation-score distribution; the two proposed
mechanisms; and algorithm-level performance followed by evidence explaining
group construction and candidate selection. Present the two mechanisms as
complementary design roles in one algorithm. Do not lead with selection-overlap
statistics or imply that the relation network learns a within-group ordering.}

The contributions are twofold. First, a hybrid PBI-based quality grouping
method combines a continuous PBI score with the binary representative-based
label used in REMO. The continuous score distinguishes solutions assigned the
same binary label when the positive training group is constructed. Second, a
dual-mode candidate selection method reduces the reliance of expensive
evaluations on a single relation score. Relation-guided candidate generation is
retained, while indicator-based reranking and exploration-oriented selection
provide two alternative criteria for selecting the final evaluation batch.

\section{Related Work}
\label{sec:related}

\subsection{Surrogate-Assisted Expensive Multi-/Many-Objective Optimization}
\label{sec:related:saea}
\TODO{Kriging/RBF-based SAEAs, their scalability limit in $M$, and why
regression on every objective becomes unattractive when $M$ grows.}

\subsection{Relation Learning for Expensive Optimization}
\label{sec:related:relation}
\TODO{REMO and the pairwise-comparison SAEA line; compare how each method forms
training groups or pairwise labels, then position the present work as a hybrid
grouping signal that combines continuous quality ranking with a binary
representative-based group label.}

\subsection{Quality Assessment and Candidate Selection}
\label{sec:related:selection}
\TODO{Indicator-based assessment (SDE, $L_p$ shape estimation, PIEA), radial
representative selection (RSEA), and infill/candidate selection rules. End with
the gap that motivates Section~\ref{sec:method}.}

\section{The Proposed HPDC-MaOEA}
\label{sec:method}

Relation-based surrogate-assisted evolutionary algorithms replace multi-output
regression by a classifier that predicts the relative quality of a pair of
solutions~\cite{ref:remo}. Two interfaces then decide what such an algorithm
can learn and how much its predictions cost: the rule that turns the evaluated
population into supervision labels, and the rule that turns the surrogate
scores of a large candidate pool into a small batch of expensive evaluations.
HPDC-MaOEA (hybrid-PBI dual-mode-candidate many-objective evolutionary
algorithm) contributes one mechanism at each interface. \emph{Hybrid PBI-based
quality grouping} (Section~\ref{sec:method:hpc}) combines a continuous PBI
score with a binary representative-based label. The score can distinguish
solutions assigned the same binary label when the positive training group is
constructed. \emph{Dual-mode candidate selection}
(Section~\ref{sec:method:candidate}) retains relation-guided candidate
generation but uses either indicator-based reranking or exploration-oriented
selection to choose the final expensive-evaluation batch. Established relation
learning, indicator assessment and environmental selection modules provide the
host algorithm; the two proposed mechanisms change how its training groups are
formed and how its surrogate scores are used.

\subsection{Overall Framework}
\label{sec:method:framework}

We consider the expensive many-objective problem
\begin{equation}
\min_{\vx\in\Omega}\;
\vf(\vx)=\bigl[f_1(\vx),\ldots,f_M(\vx)\bigr],
\qquad
\Omega\subseteq\mathbb{R}^{D},
\label{eq:problem}
\end{equation}
in which every evaluation of $\vf$ is expensive, so the number of real function
evaluations is limited to a budget $FE_{\max}$. Let $FE$ be the number of
evaluations consumed so far and $t=FE/FE_{\max}\in[0,1]$ the evaluation
progress, which both proposed mechanisms use as their only schedule variable.

HPDC-MaOEA evaluates a Latin hypercube design of $N_{\mathrm{init}}$ solutions,
stores them in an archive $\mathcal{A}_{\mathrm{rc}}$, and then repeats five
steps until the budget is exhausted. (i)~The current population $\mathcal{P}$
of size $N$ is divided into a positive group and a non-positive group by
hybrid PBI grouping, which also returns the set $\mathcal{R}$ of
representative solutions reused later as an additional mating pool.
(ii)~Ordered pairs are generated from the two groups and a three-class relation
network is trained, together with a held-out pairwise error rate $e_r$.
(iii)~An indicator value is computed for every evaluated solution and an
RBF-SVR model is fitted from decision vectors to those values.
(iv)~A relation-guided inner search accumulates a candidate pool, one of the
two selection criteria is drawn, and a batch of at most $n_{\max}$ candidates is
returned. (v)~The batch is evaluated, appended to
$\mathcal{A}_{\mathrm{rc}}$, and the next population is obtained by the
radial-grid environmental selection of RSEA~\cite{ref:rsea} over the whole
archive.

Algorithm~\ref{alg:framework} states the procedure. All expensive evaluations
after initialisation are consumed in step~(v), so both proposed mechanisms operate on already
evaluated data and surrogate predictions only, and neither changes the total
number of real evaluations. The two interfaces stay separate: grouping
decides \emph{what the classifier is asked to learn}, candidate selection
decides \emph{how its scores are spent}.

%% \begin{figure*}[t]
%%   \centering
%%   \includegraphics[width=\textwidth]{fig_framework}
%%   \caption{Overall framework of HPDC-MaOEA. \label{fig:framework}}
%% \end{figure*}
\TODO{Export the overall framework diagram as \texttt{fig\_framework.pdf} and
uncomment the figure* block above; it should show the outer loop of
Algorithm~\ref{alg:framework} with the two contributed blocks highlighted.}

\begin{algorithm}[t]
\caption{HPDC-MaOEA}
\label{alg:framework}
\begin{algorithmic}[1]
\Require problem, budget $FE_{\max}$, population size $N$
\Ensure archive $\mathcal{A}_{\mathrm{rc}}$
\State Sample $N_{\mathrm{init}}$ solutions by Latin hypercube design and
       evaluate them
\State $\mathcal{P}\gets$ initial population;
       $\mathcal{A}_{\mathrm{rc}}\gets\mathcal{P}$
\While{$FE<FE_{\max}$}
  \State $t\gets FE/FE_{\max}$
  \State $(\mathcal{C}_1,\mathcal{C}_2,\mathcal{R})\gets$
         \Call{HybridGroup}{$\mathcal{P},t$}
         \Comment{Section~\ref{sec:method:hpc}}
  \State Build ordered relation pairs from $\mathcal{C}_1,\mathcal{C}_2$ by
         \eqref{eq:pairlabel}
  \If{no pair can be formed}
    \State $\mathcal{P}\gets$
           \Call{EnvSelect}{$\mathcal{A}_{\mathrm{rc}},N$};
           \textbf{continue}
  \EndIf
  \State Train the relation surrogate and record its held-out error $e_r$
         \Comment{Section~\ref{sec:method:relation}}
  \State Train the indicator surrogate $\widehat I$ on the evaluated solutions
         \Comment{Section~\ref{sec:method:candidate:ind}}
  \State Draw the selection mode $m$ by \eqref{eq:mode}
  \State $\mathcal{S}\gets$
         \Call{DualModeSelection}{$\mathcal{P},\mathcal{R},m,t,e_r$}
         \Comment{Algorithm~\ref{alg:candidate}}
  \If{$\mathcal{S}=\emptyset$}
    \State $\mathcal{S}\gets$ at most $n_{\min}$ offspring generated from
           $\mathcal{P}\cup\mathcal{R}$
  \EndIf
  \State Truncate $\mathcal{S}$ by the remaining budget, evaluate it, and
         append it to $\mathcal{A}_{\mathrm{rc}}$
  \State $\mathcal{P}\gets$ \Call{EnvSelect}{$\mathcal{A}_{\mathrm{rc}},N$}
\EndWhile
\end{algorithmic}
\end{algorithm}

\subsection{Hybrid PBI-Based Quality Grouping}
\label{sec:method:hpc}

Relation learning first needs to divide the evaluated population into groups
from which pairwise labels can be generated. The binary PBI partition of
REMO~\cite{ref:remo} provides a clear group boundary, but all solutions assigned
the same label receive the same group-level value. It therefore cannot
distinguish two solutions of different PBI quality that lie on the same side of
the boundary. Hybrid PBI-based grouping combines a continuous PBI quality score
with the binary representative-based label. The continuous score provides
additional ranking information among solutions sharing a binary label, while
the label preserves a clear population-level partition. A stage-dependent
combination of these two signals determines the positive group used to construct
relation pairs.

Throughout this subsection $\mathcal{P}=\{(\vx_i,\vf_i)\}_{i=1}^{N}$ is the
current population and $\vz^{*}=\min_i\vf_i$ the component-wise population
minimum, taken as the ideal point.

\subsubsection{Continuous PBI Quality Score}
\label{sec:method:hpc:cont}

The continuous score assesses each solution by a PBI value measured along a
direction set $\mathcal{V}=\{\vv_j\}$. For low-dimensional or small-population
cases, $\mathcal{V}$ consists of uniform reference vectors normalised to unit
length. In the many-objective regime, the directions adapt to the current
nondominated distribution and are formed from its radial unit vectors,
\begin{equation}
\mathcal{V}=\Bigl\{\vf^{(j)}\big/\bigl\|\vf^{(j)}\bigr\|_2
\;\Bigm|\;j=1,\ldots,n_{\mathrm{ND}}\Bigr\},
\label{eq:direction}
\end{equation}
which are used directly as the adaptive direction set. Equation
\eqref{eq:direction} makes both the number and positions of these directions
follow the current nondominated distribution. Uniform vectors are used when the
adaptive construction is not activated or its fallback condition is met.

Each objective vector is associated with the direction of largest cosine
similarity, $a(i)=\arg\max_j
\vf_i^{\mathsf T}\vv_j/(\|\vf_i\|_2\|\vv_j\|_2)$, and the conventional
PBI value~\cite{ref:moead} along that direction is
\begin{equation}
\begin{aligned}
d_{1,i}&=(\vf_i-\vz^{*})^{\mathsf T}\vv_{a(i)}\big/\|\vv_{a(i)}\|_2,\\
d_{2,i}&=\bigl\|(\vf_i-\vz^{*})
        -d_{1,i}\,\vv_{a(i)}/\|\vv_{a(i)}\|_2\bigr\|_2,\\
g_i^{\mathrm{PBI}}&=d_{1,i}+\theta\,d_{2,i},
\end{aligned}
\label{eq:pbi}
\end{equation}
with penalty $\theta$. The continuous quality score is the decreasing transform
\begin{equation}
S_i=\frac{1}{1+g_i^{\mathrm{PBI}}}\in(0,1],
\label{eq:score-cont}
\end{equation}
so a larger $S_i$ indicates a smaller PBI value along the associated direction.
Unlike a class label, $S_i$ can distinguish solutions assigned the same binary
label.
Equations \eqref{eq:pbi}--\eqref{eq:score-cont} also fix the coordinate
convention of the implementation: the association is computed from the original
objective vectors, whereas the two distances are computed after translation by
$\vz^{*}$. All experiments in this paper use this single convention.

\subsubsection{Representative-Based Binary Group Label}
\label{sec:method:hpc:coarse}

The continuous score supplies a fine-grained ordering, while the binary label
adds a representative-relative quality boundary. This boundary is derived from a small set
$\mathcal{R}\subset\mathcal{P}$ of $k$ representative solutions, selected from
the current population by the radial-grid representative selection of
RSEA~\cite{ref:rsea}, which is also the rule used by REMO~\cite{ref:remo}. The
representatives are evaluated solutions that act as dynamic, population-derived
reference points. Their number increases with the objective dimension to provide
finer regional coverage as the many-objective geometry becomes richer.

Given $\mathcal{R}$, the binary group label applies the PBI partition of
REMO~\cite{ref:remo} with the representatives as reference points. Every
solution is associated with the representative $\vr_{b(i)}$ of largest cosine
similarity, the associated direction is
$\vw_{b(i)}=(\vr_{b(i)}-\vz^{*})/\|\vr_{b(i)}-\vz^{*}\|_2$, the two distances
$\hat d_{1,i},\hat d_{2,i}$ are formed as in \eqref{eq:pbi} with
$\vv_{a(i)}$ replaced by $\vw_{b(i)}$, and the binary label is
\begin{equation}
L_i=\mathbb{I}\!\left[
\frac{\hat d_{1,i}+\delta\,\hat d_{2,i}}
     {\|\vr_{b(i)}-\vz^{*}\|_2}\leq1\right]\in\{0,1\}.
\label{eq:label-bin}
\end{equation}
Normalising by $\|\vr_{b(i)}-\vz^{*}\|_2$ makes the associated representative
itself the threshold, so $L_i$ records on which side of the surface through that
representative a solution lies. Following REMO, the class-balance variable
$\delta$ is adjusted by bounded bisection so that the boundary maintains useful
occupancy on both sides across changing population states.

\subsubsection{Hybrid Score and Positive-Group Construction}
\label{sec:method:hpc:hybrid}

The two signals are combined by a weight that follows the consumption of the
expensive budget. With $t=FE/FE_{\max}$ and $\alpha_t=1-t$, the hybrid quality
score of solution $i$ is
\begin{equation}
H_i=\alpha_t S_i+(1-\alpha_t)L_i ,
\qquad
\alpha_t=1-t ,
\label{eq:fuse}
\end{equation}
and the positive group collects the highest-ranked fraction $r_g$ of the
population,
\begin{equation}
\mathcal{C}_1=\bigl\{\vx_i\;\big|\;
\operatorname{rank}_{\downarrow}(H_i)\leq\lceil r_gN\rceil\bigr\},
\qquad
\mathcal{C}_2=\mathcal{P}\setminus\mathcal{C}_1 .
\label{eq:group}
\end{equation}
With $r_g<1/2$ the non-positive group $\mathcal{C}_2$ provides a broad contrast
set, while the compact positive group concentrates pairwise supervision on the
highest-priority portion of the population.

The hybrid score connects quality assessment to relation learning through the
chain
\begin{equation*}
\begin{aligned}
&\text{quality score}\;\rightarrow\;\text{positive-group construction}\\
&\qquad\rightarrow\;\text{relation-pair construction},
\end{aligned}
\end{equation*}
so a change in the ranking matters insofar as it changes
which solutions enter $\mathcal{C}_1$ and therefore which ordered pairs the
relation model is trained on. The budget-dependent weighting connects
fine-grained quality ranking with representative-based group priority. As the
evaluation budget is consumed, the binary label receives increasing weight,
while $S_i$ refines the within-label order at each training step with $FE<FE_{\max}$.
Proposition~\ref{prop:order} characterises the resulting ordering.

\begin{proposition}[Ordering induced by the hybrid score]
\label{prop:order}
Let $S_i\in(0,1]$ as guaranteed by \eqref{eq:score-cont} whenever
$\vv_j\succeq\bm{0}$ and $\vf_i\succeq\vz^{*}$, so that
$g_i^{\mathrm{PBI}}\geq0$. Then
\begin{enumerate}[label=(\roman*),leftmargin=*,itemsep=1pt]
\item for any pair with $L_i=L_j$,
      $\;H_i-H_j=\alpha_t\,(S_i-S_j)$;
\item for any pair with $L_i=1$ and $L_j=0$,
      $\;\alpha_t\leq1/2\;\Rightarrow\;H_i>H_j$.
\end{enumerate}
\end{proposition}

\begin{proof}
(i) is immediate from \eqref{eq:fuse} because the term $(1-\alpha_t)L_i$
cancels. For (ii), if $\alpha_t=0$, then $H_i=1>0=H_j$.
For $0<\alpha_t\leq1/2$, $H_i=(1-\alpha_t)+\alpha_tS_i>1-\alpha_t$ because $S_i>0$,
while $H_j=\alpha_tS_j\leq\alpha_t$ because $S_j\leq1$; for
$\alpha_t\leq1/2$ we have $1-\alpha_t\geq\alpha_t$, hence $H_i>H_j$.
\end{proof}

Part~(i) explains why the continuous score is included: for $\alpha_t>0$,
solutions assigned the same binary label can still be distinguished when the
positive group is formed. Part~(ii) gives the schedule a concrete interpretation.
For $t<1/2$, the continuous score can affect ordering across the binary boundary;
for $1/2\leq t<1$, binary-positive solutions take priority and the continuous
score refines the order within each label. This transition combines fine-grained
ranking with a clear group priority as the budget is consumed.
Section~\ref{sec:exp:hpc} examines the resulting group composition and its
association with subsequent solution retention.

\subsection{Relation Learning on the Hybrid Groups}
\label{sec:method:relation}

HPDC-MaOEA retains the relation model of REMO~\cite{ref:remo} and supplies it
with the redesigned hybrid groups. Every ordered pair
of distinct solutions becomes one training sample whose input is the
concatenation $[\vx_i,\vx_j]\in\mathbb{R}^{2D}$ and whose label
\begin{equation}
y_{ij}=
\begin{cases}
+1, & \vx_i\in\mathcal{C}_1,\ \vx_j\in\mathcal{C}_2,\\
-1, & \vx_i\in\mathcal{C}_2,\ \vx_j\in\mathcal{C}_1,\\
0,  & \text{otherwise},
\end{cases}
\label{eq:pairlabel}
\end{equation}
with $\mathcal{C}_1,\mathcal{C}_2$ from \eqref{eq:group}, records the ordered
group relation rather than a Pareto comparison. Because $r_g<1/2$ makes the
same-group families larger than the cross-group ones, they are randomly
subsampled towards the cross-group count, so the three labels of
\eqref{eq:pairlabel} are balanced by construction rather than by a reweighted
loss.

We employ the same relation-learning architecture as REMO: min--max normalised
pair inputs, a class-balanced training/held-out split, and a feed-forward pattern
network ending in a softmax over the three labels.
Its output for a pair is written
$\bm{\pi}(\vx_a,\vx_b)=[\pi_{+1},\pi_0,\pi_{-1}]$, the predicted probabilities
that the first solution belongs to the higher group, that both belong to the
same group, and that the second belongs to the higher group. The held-out
misclassification rate $e_r$ is reused by the reliability-aware acquisition
rule in Section~\ref{sec:method:candidate:exp}.

For an unevaluated candidate $\vx$ the trained network is queried on the four
families of ordered comparisons $(\mathcal{C}_1,\vx)$, $(\vx,\mathcal{C}_1)$,
$(\mathcal{C}_2,\vx)$ and $(\vx,\mathcal{C}_2)$, that is on $2N$ pairs per
candidate. With $\mathbf{a},\mathbf{b},\mathbf{c},\mathbf{d}$ the family-wise
mean probability vectors in that order, the relation score aggregates the
evidence that $\vx$ is not inferior to the positive group and is superior to the
non-positive group, minus the opposite evidence,
\begin{equation}
R(\vx)=2\bigl[c_{-1}(\vx)+d_{+1}(\vx)-a_{+1}(\vx)-b_{-1}(\vx)\bigr]
\in[-4,4],
\label{eq:relscore}
\end{equation}
the bound following from each probability vector summing to one. $R(\vx)$ is the
aggregate group-relative preference used to order unevaluated candidates.

\subsection{Dual-Mode Candidate Selection}
\label{sec:method:candidate}

The relation surrogate can rank thousands of predicted-promising candidates per
iteration, but spending the expensive budget on the top of that single ranking
concentrates the batch on whatever preference the current surrogate happens to
hold. In particular, the number of candidates passing a fixed cutoff depends on
the distribution of $R(\vx)$, which changes with the problem, search stage and
training set. The resulting batch size, and hence the number of evaluations
spent in an iteration, can vary even though the cutoff is unchanged
(Section~\ref{sec:exp:candidate}). Moreover, $R(\vx)$ expresses only a
candidate's predicted position relative to the two training groups.

HPDC-MaOEA responds with two design decisions. First, ranking and expenditure
control are separated: quantile rules define thresholds relative to the current
score distribution, while an explicit cap $n_{\max}$ bounds the evaluation
batch. Scores determine candidate priority within this bound; the retained-set
size and remaining budget determine the feasible batch size. Second, two alternative ranking
criteria are made available --- an exploration-oriented mode
(Section~\ref{sec:method:candidate:exp}) and an indicator-oriented mode
(Section~\ref{sec:method:candidate:ind}) --- and one of them is drawn per
iteration (Section~\ref{sec:method:candidate:mix}). The two modes provide two
alternative criteria for allocating expensive evaluations while keeping the
per-iteration expenditure explicit and bounded.

Both modes rank the same kind of candidate pool, produced by a relation-guided
inner search: genetic variation is applied to the union of the current
population and the representatives $\mathcal{R}$, the relation score
\eqref{eq:relscore} ranks the offspring, the best of them are retained as
parents of the next inner generation together with $\mathcal{R}$, and the loop
stops once the cumulative number of generated candidates reaches the limit
$g_{\max}$. All candidates generated along the way are accumulated and
deduplicated into the pool $\mathcal{A}$. Mode-specific relation aggregation
shapes the inner search and produces a criterion-specific pool, while the common
pool-construction rule and batch bound preserve a consistent expenditure
interface. Section~\ref{sec:exp:candidate} examines how the two criteria change
the selected batches and the resulting search performance.

\subsubsection{Exploration-Oriented Mode}
\label{sec:method:candidate:exp}

This mode supplements the relation score with two quantities the score alone
does not carry: how sharply the classifier separates the classes around a
candidate, and how far the candidate lies from the rest of the batch in
decision space. Let $\mathcal{O}(\vx)$ collect the $2N$ ordered comparisons of
$\vx$ used in \eqref{eq:relscore}. The prediction ambiguity and the
exploration-oriented acquisition score are
\begin{equation}
\begin{aligned}
U(\vx)&=1-\frac{1}{|\mathcal{O}(\vx)|}
\sum_{(\vx_a,\vx_b)\in\mathcal{O}(\vx)}\;
\max_{y}\;\pi_y(\vx_a,\vx_b),\\
A_{\mathrm{exp}}(\vx)&=\widetilde R(\vx)+\lambda_t\,\widetilde U(\vx),
\end{aligned}
\label{eq:expscore}
\end{equation}
where $\widetilde{(\cdot)}$ denotes min--max normalisation over the pool. A
large $U(\vx)$ means the network produces flatter class distributions around
that candidate. The resulting softmax-concentration signal supplies a relative
acquisition order over the current pool. In this mode the same sharpness
$\max_y\pi_y$ additionally weights the family-wise averages entering
\eqref{eq:relscore}, so that sharper pairs contribute more to the aggregated
score while preserving the semantics of the three relation classes.

The weight of the ambiguity term is reduced according to the search stage and
the error of the model whose ambiguity is being sampled,
\begin{equation}
\lambda_t=\lambda_0(1-t)\max\!\left(0,\;1-\frac{e_r}{e_{\max}}\right),
\label{eq:lambda}
\end{equation}
with $e_r$ the held-out pairwise error of
Section~\ref{sec:method:relation} and $e_{\max}$ a fixed threshold. The
reliability factor automatically suppresses ambiguity guidance as the
held-out error approaches the threshold, while the progress factor shifts the
acquisition towards relation quality as the budget is consumed.

Candidates at or above the $q_{\mathrm{keep}}$ quantile of $A_{\mathrm{exp}}$
form the retained set
\begin{equation}
\mathcal{H}^{\mathrm{exp}}=\Bigl\{\vx\in\mathcal{A}\;\Bigm|\;
A_{\mathrm{exp}}(\vx)\geq
Q_{q_{\mathrm{keep}}}\bigl(A_{\mathrm{exp}}(\mathcal{A})\bigr)\Bigr\},
\label{eq:prefilter-exp}
\end{equation}
completed with the highest-scoring candidates if fewer than $n_{\min}$ qualify.
The batch is then assembled greedily from $\mathcal{H}^{\mathrm{exp}}$: the
first member maximises $A_{\mathrm{exp}}$, and each subsequent member maximises
\begin{equation}
\begin{split}
A_{\mathrm{batch}}(\vx\mid\mathcal{S})
&=w\,\widehat A_{\mathrm{exp}}(\vx)
+(1-w)\,\widehat d(\vx,\mathcal{S}),\\
d(\vx,\mathcal{S})&=\min_{\vz\in\mathcal{S}}\|\vx-\vz\|_2,
\end{split}
\label{eq:greedy}
\end{equation}
where $\widehat{(\cdot)}$ is renormalised over the remaining candidates at each
step and $w>1/2$. Quality remains the primary criterion, while decision-space
distance distributes the expensive batch across distinct candidate regions.

\subsubsection{Indicator-Based Reranking Mode}
\label{sec:method:candidate:ind}

The second mode retains the relation model for coarse candidate filtering and
introduces a different quality criterion only for final reranking. Following
PIEA~\cite{ref:piea}, an SDE-based fitness~\cite{ref:sde} is computed for each
evaluated solution using the generalised $L_p$ front-shape estimate and its
finite convergence fallback. Because this fitness is available only for
evaluated solutions, an RBF-SVR model approximates it in the decision space,
\begin{equation}
\widehat I(\vx)=\mathcal{M}_I(\vx;\mathcal{D}),
\qquad
\mathcal{D}=\{(\vx_i,I_i)\}_{i=1}^{N},
\label{eq:svr}
\end{equation}
and reranks the pool. Selection then follows two stages. The relation score
\eqref{eq:relscore} first retains the top fraction $q_{\mathrm{rel}}$ of
$\mathcal{A}$, after which $\widehat I$ reranks this reduced set. Candidates at
or above its $q_{\mathrm{ind}}$ quantile form $\mathcal{H}^{\mathrm{ind}}$ and
are completed with the highest-ranked candidates if fewer than $n_{\min}$
qualify. The final batch contains the best $n_b$ members under $\widehat I$.
Thus, the contribution of this mode is the two-stage use of relation-based
filtering followed by indicator-based reranking.

\subsubsection{Probabilistic Mode Switching}
\label{sec:method:candidate:mix}

Let $m\in\{\mathrm{exp},\mathrm{ind}\}$ denote the mode of the current
iteration. With $u\sim U[0,1)$,
\begin{equation}
m=
\begin{cases}
\mathrm{ind}, & \text{if the indicator surrogate is available}\\
              & \text{and } u<p_{\mathrm{mix}},\\[2pt]
\mathrm{exp}, & \text{otherwise},
\end{cases}
\label{eq:mode}
\end{equation}
with $p_{\mathrm{mix}}=0.5$. Stochastic switching assigns one intact criterion
to each iteration, preserving the distinct selection semantics of the two modes
and giving them equal selection probability when the indicator surrogate is
available, without requiring their scores to be combined into a single ranking.

Whichever mode is drawn, the number of expensive evaluations is
\begin{equation}
n_b=\min\bigl(n_{\max},\,|\mathcal{H}^{m}|\bigr),
\label{eq:batch}
\end{equation}
truncated further by the remaining budget $FE_{\max}-FE$. Before this cap is
applied, the retained set is completed to $n_{\min}$ whenever the pool permits.
Together, the quantile rules and \eqref{eq:batch} convert surrogate scores into
a stable rank-and-budget interface. A direct variation step supplies candidates
when the inner pool is empty. Algorithm~\ref{alg:candidate} summarises the
module.

\begin{algorithm}[t]
\caption{Dual-mode candidate selection}
\label{alg:candidate}
\begin{algorithmic}[1]
\Require population $\mathcal{P}$, representatives $\mathcal{R}$, relation and
         indicator surrogates, mode $m$, progress $t$, held-out error $e_r$
\Ensure candidate batch $\mathcal{S}$
\State $\mathcal{A}\gets\emptyset$;
       $\mathcal{Q}\gets$ offspring of $\mathcal{P}\cup\mathcal{R}$
\While{$|\mathcal{A}|<g_{\max}$ and $\mathcal{Q}\neq\emptyset$}
  \State $\mathcal{A}\gets\mathcal{A}\cup\mathcal{Q}$
  \State Score $\mathcal{Q}$ by the relation surrogate, \eqref{eq:relscore},
         and keep the best of them as parents
  \State $\mathcal{Q}\gets$ offspring of the parents and $\mathcal{R}$
\EndWhile
\State Remove duplicates from $\mathcal{A}$
\If{$m=\mathrm{ind}$}
  \State Retain the top fraction $q_{\mathrm{rel}}$ of $\mathcal{A}$ by
         \eqref{eq:relscore}
  \State Rerank the retained subset by $\widehat I$ of \eqref{eq:svr} and keep
         its top quantile as $\mathcal{H}^{\mathrm{ind}}$
  \State $\mathcal{S}\gets$ the $n_b$ best members of
         $\mathcal{H}^{\mathrm{ind}}$ under $\widehat I$
\Else
  \State Compute $A_{\mathrm{exp}}$ by \eqref{eq:expscore}--\eqref{eq:lambda}
         and form $\mathcal{H}^{\mathrm{exp}}$ by \eqref{eq:prefilter-exp}
  \State $\mathcal{S}\gets$ greedy batch of $n_b$ members of
         $\mathcal{H}^{\mathrm{exp}}$ by \eqref{eq:greedy}
\EndIf
\State Bound $n_b$ by \eqref{eq:batch} and by the remaining budget
\State \Return $\mathcal{S}$
\end{algorithmic}
\end{algorithm}

\subsection{Computational Complexity}
\label{sec:method:complexity}

Let $N$ be the population size, $N_v=O(N)$ the number of directions, $N_c$ the
size of the deduplicated candidate pool, $D$ the number of decision variables
and $M$ the number of objectives.

In the grouping module, association and the continuous PBI computation
\eqref{eq:pbi}--\eqref{eq:score-cont} cost $O(NN_vM)$, the direction set
\eqref{eq:direction} costs one normalisation of the nondominated vectors, the
bounded bisection of \eqref{eq:label-bin} costs $O(B_\delta NkM)$ with
$B_\delta$ fixed by the bracket and the tolerance, and representative selection
adds nondominated sorting and radial-projection distances. With $N_v=O(N)$ and
$k=O(N)$ the module is dominated by $O(N^2M)$, and generating the relation pairs
adds $O(N^2D)$.

In the candidate module, scoring one candidate requires $2N$ network inputs of
dimension $2D$, so the pool costs $O(N_cND)$ in data construction and
prediction, which dominates the module. The indicator fitness costs $O(N^2M)$,
the indicator-oriented mode adds one SVR fit on the $N$ evaluated solutions plus
predictions on the screened subset, and the greedy construction \eqref{eq:greedy}
costs $O(n_{\max}N_cD)$. The per-iteration cost is therefore $O(N_cND+N^2M)$.
Neither module consumes expensive evaluations, so the total number of real
evaluations remains $FE_{\max}$ regardless of the mode sequence; the wall-clock
share of the two modules is reported in Section~\ref{sec:exp:sensitivity}.


\section{Experimental Studies}
\label{sec:exp}

\subsection{Experimental Settings}
\label{sec:exp:settings}
\TODO{Specify benchmarks, objective counts, decision dimensions, population size,
evaluation budgets, independent runs, reported indicators, statistical procedures
and platform from the retained result sets. Emphasise the $M=10,20$ comparisons
and use the verified lower-objective settings as additional coverage. Distinguish
the FE=300 candidate probe from the FE=500 grouping study. Keep the main text
focused on contribution-defining controls; place fixed implementation settings
and result-version identifiers in the reproducibility material.}

\subsection{Comparison with State-of-the-Art Algorithms}
\label{sec:exp:comparison}
\TODO{Lead with the verified IGD results in the 10- and 20-objective settings,
using the full benchmark set and all eligible baselines present in each result
set. Use a common baseline subset for comparisons across objective counts, and
keep changes in decision dimension and budget explicit. Report problemwise
effects alongside win/tie/loss counts, including WFG3. Sections~\ref{sec:exp:hpc}
and~\ref{sec:exp:candidate} then explain the behaviour of the two design components.}

\subsection{Ablation Study}
\label{sec:exp:ablation}
\TODO{Organise available ablations around the two proposed components. Use
matched host algorithms and representative counts for component attribution,
and identify any comparison that also changes $k$. Keep missing configurations
explicit in the tables. For the candidate module, compare relation top-$6$,
exploration-only, indicator-only and dual-mode selection. Separate component
removal from a test of interaction; report a joint gain only when the matched
comparisons support it.}

\subsection{Analysis of Hybrid PBI-Based Quality Grouping}
\label{sec:exp:hpc}
\TODO{Explain how continuous ranking and the binary boundary shape
the positive group. The equal-quota comparison uses score\_hybrid, score\_v and
anchor\_margin; report the native binary-label group separately. Production
Hybrid combines score\_v with label\_dyn, not with anchor\_margin.
For the original five-problem, 250-run study, report population\_final
Precision@25\% of 0.28698, 0.27748 and 0.27617 for Hybrid, direction score and
Anchor margin, respectively, with stagewise results. Explain the distinct
information sources using the 33/40 cells with bidirectional unique future true
positives; also state that strict simultaneous fusion superiority passed in
1/40 cells. These results describe future-retention association along the
recorded Hybrid trajectories. Use the ordering proposition for the schedule's
design rationale; include alternative schedules as performance evidence only
when such comparisons are available. Place Jaccard, source-composition details
and full cellwise tests in the Supplementary Information.}

\subsection{Analysis of Dual-Mode Candidate Selection}
\label{sec:exp:candidate}
\TODO{Lead with the two candidate-selection roles: exploration-oriented batch
dispersion and indicator-based reranking within a bounded evaluation batch.
Use the FE=300 probe on four 20-objective problems, ten runs per arm. Identify
V0 as the REMO selection rule transplanted into the probe host, V1 as accumulated
pool relation top-$6$, V2 as exploration-only, V3 as indicator-only and V4 as
the full dual-mode strategy. Compare V4 with V0 as an overall selection-policy
comparison and with V1 to assess the additional criteria under the same batch
bound. The conclusion report gives paired IGD p-values of 0.008 and 0.375,
respectively; verify the statistical family and correction before final reporting.
Present late-stage within-pool greedy-reference gain, batch spread and retention
together: V4 versus V1 has reported mean gain ratios of about 0.105 versus 0.086
and retention rates of 0.770 versus 0.706. Retain the problemwise endpoint
results, including DTLZ7, so retention is interpreted alongside final quality.
Add fixed-cutoff audit numbers or same-state counterfactual results only after
their completed result tables have been matched to the manuscript version.}

\subsection{Robustness to Core Controls}
\label{sec:exp:sensitivity}
\TODO{Vary a compact set of contribution-defining controls and report the
stability of quality, batch size and budget use. Keep all inherited and
implementation-level settings fixed. Also report the wall-clock share of the two
contributed modules.}

\section{Conclusion}
\label{sec:conclusion}
\TODO{Reinforce the two design contributions: hybrid PBI-based grouping combines
fine-grained ordering with representative-based group priority, and dual-mode
candidate selection offers exploration and indicator criteria within an
explicit evaluation-batch bound. State the strongest verified algorithm-level
result and its benchmark scope. Close by connecting the group-composition and
candidate-batch findings to these two design roles, keeping empirical performance
claims at the level supported by their matched comparisons.}

\begin{thebibliography}{9}

\bibitem{ref:moead}
Q.~Zhang and H.~Li,
``MOEA/D: A multiobjective evolutionary algorithm based on decomposition,''
\emph{IEEE Trans. Evol. Comput.}, vol.~11, no.~6, pp.~712--731, 2007.

\bibitem{ref:remo}
\TODO{REMO -- relation-model-based expensive multi-objective optimization.}

\bibitem{ref:rsea}
\TODO{RSEA -- radar-grid based selection evolutionary algorithm.}

\bibitem{ref:piea}
Y.~Li, W.~Li, S.~Li, and Y.~Zhao,
``PIEA,'' \emph{Information Sciences}, 2024.
\TODO{complete volume and pages.}

\bibitem{ref:sde}
M.~Li, S.~Yang, and X.~Liu,
``Shift-based density estimation for Pareto-based algorithms in
many-objective optimization,''
\emph{IEEE Trans. Evol. Comput.}, vol.~18, no.~3, pp.~348--365, 2014.

\bibitem{ref:pcsaea}
\TODO{PC-SAEA -- pairwise-comparison-based surrogate-assisted EA,
Swarm Evol. Comput., 2023, doi:10.1016/j.swevo.2023.101323.}

\bibitem{ref:r2aea}
\TODO{R2AEA -- two-stage relation-based expensive optimization.}

\bibitem{ref:platemo}
Y.~Tian, R.~Cheng, X.~Zhang, and Y.~Jin,
``PlatEMO: A MATLAB platform for evolutionary multi-objective optimization,''
\emph{IEEE Comput. Intell. Mag.}, vol.~12, no.~4, pp.~73--87, 2017.

\end{thebibliography}

\end{document}
